
\documentclass[12pt]{article}

% Layout.
\usepackage[top=1.2in, bottom=0.75in, left=1in, right=1in, headheight=1.0in, headsep=0pt]{geometry}

% Fonts.
\usepackage{mathptmx}
\usepackage[scaled=0.86]{helvet}
\renewcommand{\emph}[1]{\textsf{\textbf{#1}}}

% TiKZ.
\usepackage{tikz, pgfplots}
\usetikzlibrary{calc}
\pgfplotsset{my style/.append style={axis x line=middle, axis y line=middle, xlabel={$x$}, ylabel={$y$}}}
\pgfplotsset{compat=1.16}

% Misc packages.
\usepackage{amsmath,amssymb,latexsym,bm,array,multicol,enumitem}
\usepackage{graphicx}
\usepackage{xcolor}
\usepackage[most]{tcolorbox}
\usepgfplotslibrary{fillbetween}
\definecolor{solbg}{RGB}{225,240,252}
\definecolor{solframe}{RGB}{70,130,180}
\newtcolorbox{solution}[1][]{%
  breakable, enhanced,
  colback=solbg, colframe=solframe, boxrule=0.7pt, arc=3pt,
  left=6pt, right=6pt, top=5pt, bottom=5pt,
  before skip=6pt, after skip=6pt,
  fonttitle=\sffamily\bfseries\small, coltitle=white, title={Solution},
  #1}

% Commands to set various header/footer components.
\makeatletter
\def\doctitle#1{\gdef\@doctitle{#1}}
\doctitle{Use {\tt\textbackslash doctitle\{MY LABEL\}}.}
\def\docdate#1{\gdef\@docdate{#1}}
\docdate{Use {\tt\textbackslash docdate\{MY DATE\}}.}
\def\doccourse#1{\gdef\@doccourse{#1}}
\let\@doccourse\@empty
\def\docscoring#1{\gdef\@docscoring{#1}}
\let\@docscoring\@empty
\def\docversion#1{\gdef\@docversion{#1}}
\let\@docversion\@empty
\makeatother

% Headers and footers layout.
\makeatletter
\usepackage{fancyhdr}
\pagestyle{fancy}
\fancyhf{} % Clears all headers/footers.
\lhead{\emph{\@doctitle\hfill\@docdate} \medskip
\ifnum \value{page} > 1\relax\else\\
\emph{Name: \rule{3.5in}{1pt}\hfill \@docscoring}
\fi}

\rfoot{\emph{\@docversion}}
\lfoot{\emph{\@doccourse}}
\cfoot{\emph{\thepage}}
\renewcommand{\headrulewidth}{0pt}%
\makeatother

% Paragraph spacing
\parindent 0pt
\parskip 6pt plus 1pt

% A problem is a section-like command. Use \problem{5} for a problem worth 5 points.
\newcounter{probcount}
\newcounter{subprobcount}
\setcounter{probcount}{0}

\newcommand{\problem}[1]{%
\par
\addvspace{4pt}%
\setcounter{subprobcount}{0}%
\stepcounter{probcount}%
\makebox[0pt][r]{\emph{\arabic{probcount}.}\hskip1ex}\emph{[#1 points]}\hskip1ex}

\newcommand{\thesubproblem}{\emph{\alph{subprobcount}.}}

% Subproblems are an enumerate-like environment with a consistent
% numbering scheme.  Use \begin{subproblems}\item...\item...\end{subproblems}
\newenvironment{subproblems}{%
\begin{enumerate}[itemindent=8pt,leftmargin=0pt]%
\setcounter{enumi}{\value{subprobcount}}%
\renewcommand{\labelenumi}{\emph{\textsl{\alph{enumi})}} \,}%
}%
{%
\setcounter{subprobcount}{\value{enumi}}%
\end{enumerate}%
}

% Blanks for answers in normal and math mode.
\newcommand{\blank}[1]{\rule{#1}{0.75pt}}
\newcommand{\mblank}[1]{\underline{\hspace{#1}}}
\def\emptybox(#1,#2){\framebox{\parbox[c][#2]{#1}{\rule{0pt}{0pt}}}}

% Misc.
\renewcommand{\d}{\displaystyle}
\newcommand{\ds}{\displaystyle}
\newcommand{\threeopts}{{\small \hspace{-6mm} $\begin{matrix} \text{\textsc{converges}} \\ \text{\textsc{absolutely}} \end{matrix}$ \qquad\qquad $\begin{matrix} \text{\textsc{converges}} \\ \text{\textsc{conditionally}} \end{matrix}$ \qquad\qquad \textsc{diverges}} \bigskip}


\newcommand{\ba}{\mathbf{a}}
\newcommand{\bb}{\mathbf{b}}
\newcommand{\bc}{\mathbf{c}}
\newcommand{\bi}{\mathbf{i}}
\newcommand{\bj}{\mathbf{j}}
\newcommand{\bk}{\mathbf{k}}
\newcommand{\bu}{\mathbf{u}}
\newcommand{\bv}{\mathbf{v}}
\newcommand{\bw}{\mathbf{w}}

\newcommand{\ip}[2]{\mathrm{\left<#1,#2\right>}}
\newcommand{\vv}[2]{\mathrm{\left<#1,#2\right>}}
\newcommand{\vvv}[3]{\mathrm{\left<#1,#2,#3\right>}}


\doctitle{Math 252: Quiz 6 Solutions}
\docdate{22 October, 2026}
\doccourse{}
\docversion{Solutions}
\docscoring{\fbox{{\LARGE \strut}\hspace{0.8in} / 25}}

\begin{document}
30 minutes maximum.  No aids (book, calculator, etc.) are permitted.
Show all work and use proper notation for full credit.  Answers should
be in reasonably-simplified form.  25 points possible.

\begin{enumerate}

\item Consider the following series $\ds \sum_{k=1}^{\infty} \left( \sqrt{k} - \sqrt{k+1} \right)$.
  \begin{enumerate}
  \item (2 points) Write out the first four partial sums.
    \begin{solution}
      $\begin{aligned}
        S_1 &= 1 - \sqrt2 \\
        S_2 &= \big(1 - \sqrt2\big) + \big(\sqrt2 - \sqrt3\big) = 1 - \sqrt3 \\
        S_3 &= \big(1 - \sqrt2\big) + \big(\sqrt2 - \sqrt3\big) + \big(\sqrt3 - 2\big) = 1 - 2 = -1 \\
        S_4 &= \big(1 - \sqrt2\big) + \big(\sqrt2 - \sqrt3\big) + \big(\sqrt3 - 2\big) + \big(2 - \sqrt5\big) = 1 - \sqrt5
      \end{aligned}$
    \end{solution}
    \vfill
    
  \item (3 points) Write a formula for the $n$th partial sum.
    \begin{solution}
      This is a telescoping series.
      \[
        S_n = \big(\sqrt1 - \sqrt2\big) + \big(\sqrt2 - \sqrt3\big) + \cdots + \big(\sqrt{n} - \sqrt{n+1}\big)
        = \boxed{1 - \sqrt{n+1}}
      \]
    \end{solution}
    \vfill
    
  \item (5 points) Determine whether the series converges or diverges. If the
    series converges, find its value.
    \begin{solution}
      \[
        \lim_{n\to\infty} S_n = \lim_{n\to\infty} \left( 1 - \sqrt{n+1} \right) = -\infty
      \]
      The partial sums diverge, so the series \boxed{\text{diverges}}.
    \end{solution}
    \vfill
  \end{enumerate}

\item (5 points) Determine whether the following series converges or
  diverges. If it converges, find its value.
  \[
    \sum_{n=1}^{\infty} 500 \left( \frac{1}{5^n} \right)
  \]
  \begin{solution}
    Shifting the index to start at $n = 0$,
    \[
      \sum_{n=1}^{\infty} 500 \left( \frac{1}{5^n} \right)
      = \sum_{n=0}^{\infty} 500 \left( \frac{1}{5^{n+1}} \right)
      = \sum_{n=0}^{\infty} \frac{500}{5} \left( \frac15 \right)^{n}
      = \sum_{n=0}^{\infty} 100 \left( \frac15 \right)^{n}.
    \]
    This is a geometric series with $a = 100$ and $r = \frac15$. Since $|r| = \frac15 < 1$,
    it converges, and
    \[
      \sum_{n=0}^{\infty} 100 \left( \frac15 \right)^{n} = \frac{a}{1-r} = \frac{100}{1 - \frac15}
      = \frac{100}{\frac45} = \boxed{125}
    \]
  \end{solution}
  \vfill

  \newpage
  
\item For each of the following series, determine whether the series converges or diverges.
  \begin{enumerate}
  \item (4 points) $\ds \sum_{n=5}^{\infty} 2^{(1/n)}$
    \begin{solution}
      \[
        \lim_{n\to\infty} 2^{1/n} = 2^0 = 1 \neq 0,
      \]
      so by the divergence test the series \boxed{\text{diverges}}.
    \end{solution}
    \vfill
    
  \item (4 points) $\ds \sum_{n=1}^{\infty} \frac{n}{1+n^2}$
    \begin{solution}
      Use the integral test with $f(x) = \dfrac{x}{1+x^2}$, so that $f(n) = \dfrac{n}{1+n^2}$.
      On $[1,\infty)$, $f$ is positive and continuous, and it is decreasing since
      $f'(x) = \dfrac{1-x^2}{(1+x^2)^2} \le 0$ for $x \ge 1$.
      \begin{align*}
        \int_1^\infty \frac{x}{1+x^2}\,dx
          &= \lim_{t\to\infty} \int_1^t \frac{x}{1+x^2}\,dx
           = \lim_{t\to\infty} \Big[ \tfrac12 \ln\!\big(1+x^2\big) \Big]_1^t \\[4pt]
          &= \lim_{t\to\infty} \left( \tfrac12 \ln\!\big(1+t^2\big) - \tfrac12 \ln 2 \right) = \infty
      \end{align*}
      The integral diverges, so by the integral test, the series \boxed{\text{diverges}}.
    \end{solution}
    \vfill
    \vfill
    
  \item (2 points) $\ds \sum_{n=1}^{\infty} \frac{1}{n^4}$
    \begin{solution}
      This is a $p$-series with $p = 4 > 1$, so it \boxed{\text{converges}}.
    \end{solution}
    \vfill
  \end{enumerate}

\end{enumerate}

\end{document}

%-------------------------------------------------------------------------------------------------------------------------------------------------------------------------------------------------------------------

%%% Local Variables:
%%% mode: latex
%%% TeX-master: t
%%% End:
